\newcommand{\K}{\mathbb{K}}
\newcommand{\F}{\mathbb{F}}
\newcommand{\Fp}{\mathbb{F}_p}
\newcommand{\Fq}{\mathbb{F}_q}
\newcommand{\R}{\mathbb{R}}
\newcommand{\Z}{\mathbb{Z}}
\newcommand{\V}{\mathbb{V}}

\newcommand{\idmat}{\ensuremath{\mathbf{I}}}
\newcommand{\ideal}[1]{\ensuremath{\langle #1 \rangle}}
\newcommand{\spoly}{\ensuremath{\mathcal{S}}}

\newcommand{\bigO}{\mathcal{O}}
\renewcommand{\deg}{\ensuremath{\mathsf{deg}}}
\newcommand{\vars}[1]{\mathbf{#1}}

% \newcommand{\ubar}{\underset{\bar{}}}
\newcommand{\ubar}[1]{\underline{#1}}
% \newcommand\ubar[1]{\stackunder[1.2pt]{$#1$}{\rule{.8ex}{.075ex}}}
% \newcommand{\Res}{\mathsf{Res}}
\newcommand{\Lin}{\mathsf{Lin}}
\newcommand{\vin}{\mathsf{gens}}
\newcommand{\supp}{\mathsf{supp}}
\newcommand{\coeff}{\mathsf{coeff}}
\newcommand{\exponents}{\mathsf{degs}}
\newcommand{\rref}{\mathsf{rref}}
\newcommand{\lm}{\mathsf{LM}}
\newcommand{\lc}{\mathsf{LC}}
\newcommand{\lt}{\mathsf{LT}}
\newcommand{\initideal}{\mathrm{In}}
\newcommand{\wt}{\mathsf{wt}}
% \newcommand{\dim}{\mathrm{dim}}
% \newcommand{\rank}{\mathsf{Rank}}
\newcommand{\boldparagraph}[1]{\vspace{0.5cm}
\textbf{#1} }
\newcommand{\variables}[2]{x_#1, \dots, x_#2}

\newcommand\nth{\textsuperscript{th}\xspace}
\newcommand{\etal}{\textit{et al. }}


\newcommand{\ufnperm}{\ensuremath{\text{UFN}^\pi}}
\newcommand{\bfnperm}{\ensuremath{\text{Feistel-II}^\pi}}
\newcommand{\typethreeperm}{\ensuremath{\text{Feistel-III}^\pi}}
\newcommand{\mrfnperm}{\ensuremath{\text{MRFN}^\pi}}
\newcommand{\ufnerf}{\ensuremath{\text{UFN}_{\text{erf}}}}
\newcommand{\ufncrf}{\ensuremath{\text{UFN}_{\text{crf}}}}
\newcommand{\bfnnyb}{\ensuremath{\text{Feistel-II}}}
\newcommand{\bfntwo}{\ensuremath{\text{Feistel-II}}}
\newcommand{\bfnthree}{\ensuremath{\text{Feistel-III}}}
\newcommand{\mrfn}{\ensuremath{\text{MRFN}}}


% for editorial purpose
%
\newcommand{\ar}[1]{\textcolor{blue}{AR: #1}}
\newcommand{\luc}[1]{\textcolor{red}{LC: #1}}
\newcommand{\mlf}[1]{\textcolor{green}{MX: #1}}
% \newlist{todolist}{itemize}{2}
% \setlist[todolist]{label=$\square$}


\newcommand{\YES}{\checkmark}
\newcommand{\NO}{\texttimes}
\newcommand{\used}{YES}
\newcommand{\notused}{NO}
\newcommand{\slwe}{search-LWE} 
\newcommand{\belwe}{BinaryError-LWE}
\newcommand{\bslwe}{BinarySecret-LWE}
\newcommand{\LT}{\ensuremath{\mathrm{LT}}}
\newcommand{\LM}{\ensuremath{\mathrm{LM}}}
\newcommand{\dreg}{\ensuremath{d_{\textsf{reg}}}}


% \tikzset{XOR/.style={draw,circle,thick,inner sep=1.2mm,append after command={
%         [shorten >=\pgflinewidth, shorten <=\pgflinewidth,]
%         (\tikzlastnode.north) edge[thick] (\tikzlastnode.south)
%         (\tikzlastnode.east) edge[thick] (\tikzlastnode.west)
%         }
%     }
%   }%


%% names

\newcommand{\ipos}{\ensuremath{\mathtt{IPoS}}}
\newcommand{\matrixdef}[1]{\ensuremath{#1}}

%% For polynomial division

%% for multivariate polynomial division
% \begin{array}{rll}
%      a_1\colon  & \multicolumn{1}{l}{x+y}\\
%      a_2\colon  & \multicolumn{1}{l}{1}           & \ r\\ \cline{1-2}
%   xy + 1\PhantC & \multirow{2}*{$\bigg) \ x^2y + xy^2+y^2 $} \\\cline{3-3}
%  y^2 + 1\PhantC &\\
%     %
%     &\PhantSQ           x^2y -     x                                         \\\CMidRule[3.0ex][9.0ex]{2-2}
%     &\PhantSQ \hphantom{x^2y +{}}  xy^2 +    x + y^2                         \\
%     &\PhantSQ \hphantom{x^2y +{}}  xy^2 -    y                               \\\CMidRule[9.0ex][5.0ex]{2-2}
%     &\PhantSQ \hphantom{x^2y +     xy^2 +{}} x +    y^2 +     y              \\\CMidRule[16.0ex][5.0ex]{2-2}
%     &\PhantSQ \hphantom{x^2y +     xy^2 +    x +{}} y^2 +     y              \\
%     &\PhantSQ \hphantom{x^2y +     xy^2 +    x +{}} y^2 -     1              \\\CMidRule[20.0ex][5.0ex]{2-2}
%     &\PhantSQ \hphantom{x^2y +     xy^2 +    x +    y^2 +{}}  y + 1          \\\CMidRule[25.0ex][1.0ex]{2-2}
%     &\PhantSQ \hphantom{x^2y +     xy^2 +    x +    y^2 + y +{}}  1 &\to x+y \\\CMidRule[25.0ex][1.0ex]{2-2}
%     &\PhantSQ \hphantom{x^2y +     xy^2 +    x +    y^2 + y +{}}  0 &\to x+y+1
%   \\&\PhantSQ\color{red}x^2y +     xy^2 +    x +    y^2 + y + 1
%    % Last line above should be removed -- used for alignment purposes only.
% \end{array}